%==============================================================================
%  Presentación de seminario
%  "El efecto de la forma de la parcela sobre el precio del suelo urbano"
%  Glumac, Herrera-Gómez, Licheron & Francke
%  Beamer — compilar en Overleaf con pdfLaTeX (o XeLaTeX)
%==============================================================================
\documentclass[11pt,aspectratio=169]{beamer}

%------------------------------ Tema y colores --------------------------------
\usetheme{Madrid}
\usecolortheme{default}

\definecolor{azulprincipal}{RGB}{29,53,87}
\definecolor{azulsecundario}{RGB}{69,123,157}
\definecolor{verdecompacto}{RGB}{43,124,68}
\definecolor{naranjaextruido}{RGB}{200,80,40}

\setbeamercolor{structure}{fg=azulprincipal}
\setbeamercolor{frametitle}{bg=azulprincipal,fg=white}
\setbeamercolor{title}{fg=azulprincipal}
\setbeamercolor{block title}{bg=azulsecundario,fg=white}
\setbeamercolor{block body}{bg=azulsecundario!10}
\setbeamertemplate{navigation symbols}{}
\setbeamertemplate{caption}[numbered]

%------------------------------ Idioma y fuentes ------------------------------
\usepackage[utf8]{inputenc}
\usepackage[T1]{fontenc}
\usepackage[spanish,es-tabla]{babel}
% Evita que babel-spanish redefina caracteres activos que rompen tablas/URLs
\spanishdecimal{,}
\usepackage{lmodern}

%------------------------------ Paquetes útiles -------------------------------
\usepackage{booktabs}
\usepackage{array}
\usepackage{amsmath,amssymb}
\usepackage{graphicx}
\usepackage{tikz}
\usepackage{xcolor}
\usepackage{colortbl}
\usepackage{multirow}
\usetikzlibrary{shapes.geometric,positioning}

%------------------------------ Comandos auxiliares ---------------------------
\newcommand{\Gindex}{\textit{G-index}\xspace}
\usepackage{xspace}

% Pie de página compacto
\setbeamertemplate{footline}{%
  \leavevmode%
  \hbox{%
  \begin{beamercolorbox}[wd=.40\paperwidth,ht=2.25ex,dp=1ex,center]{author in head/foot}%
    \usebeamerfont{author in head/foot}\insertshortauthor
  \end{beamercolorbox}%
  \begin{beamercolorbox}[wd=.45\paperwidth,ht=2.25ex,dp=1ex,center]{title in head/foot}%
    \usebeamerfont{title in head/foot}\insertshorttitle
  \end{beamercolorbox}%
  \begin{beamercolorbox}[wd=.15\paperwidth,ht=2.25ex,dp=1ex,right]{date in head/foot}%
    \usebeamerfont{date in head/foot}\insertframenumber{}/\inserttotalframenumber\hspace*{2ex}
  \end{beamercolorbox}}%
}

%==============================================================================
%  Datos del título
%==============================================================================
\title[Forma de la parcela y precio del suelo]{\color{white}{El efecto de la forma de la parcela\\sobre el precio del suelo urbano}}
\subtitle{\color{white}{Primera evidencia empírica del \Gindex}}
\author[Glumac, Herrera-Gómez, Licheron \& Francke]{Brano Glumac \and Marcos Herrera-Gómez \and \\ Julien Licheron \and Marc Francke}
\institute[]{\small Henley Business School \, $\cdot$ \, CONICET--CIANECO, UNRC \\ LISER \, $\cdot$ \, University of Amsterdam}
\date{Seminario de investigación \\ \today}

%==============================================================================
\begin{document}

%------------------------------ Portada ---------------------------------------
\begin{frame}[plain]
  \titlepage
\end{frame}

%------------------------------ Agenda ----------------------------------------
\begin{frame}{Agenda}
  \tableofcontents
\end{frame}

%==============================================================================
\section{Motivación}
%==============================================================================
\begin{frame}{Motivación: ¿la forma de la parcela tiene precio?}
  \begin{itemize}
    \item El precio del suelo urbano depende de un \alert{conjunto de atributos}: tiempo, regulación, localización y características físicas.
    \item El \textbf{tamaño} de la parcela es un determinante muy estudiado del valor del suelo (Colwell \& Munneke, 1997).
    \item La \textbf{geometría}, en cambio, afecta \emph{cómo} puede usarse el suelo, no solo \emph{cuánto} suelo hay.
    \item Pese a su relevancia intuitiva, su efecto marginal sobre el precio \alert{nunca había sido cuantificado directamente}.
  \end{itemize}
  \vspace{0.4em}
  \begin{block}{Pregunta central}
    ¿En qué medida la forma de una parcela individual se capitaliza en su precio de transacción?
  \end{block}
\end{frame}

\begin{frame}{La forma importa en distintas escalas espaciales}
  \begin{columns}[T]
    \begin{column}{0.55\textwidth}
      \begin{itemize}
        \item \textbf{Escala edificio} (Lindenthal, 2020): la forma y la fachada influyen en costos y valor.
        \item \textbf{Escala ciudad} (Harari, 2020): la geometría urbana afecta productividad, transporte y energía.
        \item \textbf{Escala intermedia} $\rightarrow$ la \alert{parcela individual}, unidad operativa de las decisiones de desarrollo.
      \end{itemize}
      \vspace{0.6em}
      Si la geometría restringe lo que se puede construir, el mercado debería \alert{ponerle precio} a esa restricción.
    \end{column}
    \begin{column}{0.45\textwidth}
      \centering
      \begin{tikzpicture}[scale=0.9]
        % Edificio
        \node at (0,2.4) {\footnotesize Edificio};
        \draw[fill=azulsecundario!30] (-0.4,1.4) rectangle (0.4,2.1);
        % Ciudad
        \node at (0,1.0) {\footnotesize Ciudad};
        \draw[fill=azulsecundario!30] (-0.9,-0.1) -- (0.9,0.2) -- (0.6,0.7) -- (-0.7,0.6) -- cycle;
        % Parcela
        \node at (0,-0.7) {\footnotesize \textbf{Parcela}};
        \draw[fill=naranjaextruido!30,thick] (-0.8,-1.6) rectangle (0.8,-1.1);
        \node[align=center] at (0,-2.1) {\scriptsize unidad de\\[-2pt]\scriptsize desarrollo};
      \end{tikzpicture}
    \end{column}
  \end{columns}
\end{frame}

\begin{frame}{¿Por qué la forma importa económicamente?}
  Las parcelas \alert{extruidas} (irregulares o alargadas) son más difíciles de usar, subdividir o ensamblar. Canales de impacto:
  \vspace{0.5em}
  \begin{itemize}
    \item \textbf{Envolvente edificable}: los retiros obligatorios desperdician más perímetro $\Rightarrow$ menor superficie útil.
    \item \textbf{Flexibilidad arquitectónica}: las formas raras exigen diseños a medida y comprometen la densidad.
    \item \textbf{Infraestructura y acceso}: calles, accesos y servicios son más costosos en lotes irregulares.
    \item \textbf{Subdivisión y zonificación}: pueden violar requisitos de frente o fondo mínimos.
    \item \textbf{Riesgo y reventa}: menos comparables $\Rightarrow$ los tasadores aplican descuentos.
  \end{itemize}
  \vspace{0.4em}
  \begin{center}
    \itshape Hipótesis: a igualdad de condiciones, las parcelas más extruidas se transan con descuento.
  \end{center}
\end{frame}

%==============================================================================
\section{Medición de la forma: el \texorpdfstring{\Gindex}{G-index}}
%==============================================================================
\begin{frame}{Del índice de Gravelius al \Gindex}
  \begin{columns}[T]
    \begin{column}{0.5\textwidth}
      \textbf{Índice de Gravelius (1914)}
      \[ G_n=\frac{P_n}{2\sqrt{\pi A_n}}=0.28\,\frac{P_n}{\sqrt{A_n}} \]
      \begin{itemize}\footnotesize
        \item $=1$ para un círculo (máxima compacidad).
        \item Adimensional e invariante a escala.
        \item Referencia poco realista para el suelo urbano.
      \end{itemize}
    \end{column}
    \begin{column}{0.5\textwidth}
      \textbf{\Gindex (propuesto)}
      \[ \text{\Gindex}_n=\frac{P_n}{4\sqrt{A_n}} \]
      \begin{itemize}\footnotesize
        \item $=1$ para un \alert{cuadrado} (referencia de desarrollo eficiente).
        \item $>1$ $\Rightarrow$ mayor extrudidad.
        \item Transformación lineal: $\text{\Gindex}_n=\tfrac{\sqrt\pi}{2}G_n\approx 0.886\,G_n$.
      \end{itemize}
    \end{column}
  \end{columns}
  \vspace{0.6em}
  \begin{block}{Idea clave}
    Misma información geométrica, \alert{mejor interpretación económica}: anclar la compacidad a un cuadrado se alinea con la planificación urbana ortogonal.
  \end{block}
\end{frame}

\begin{frame}{El \Gindex en cuatro formas de igual área (100 u.)}
  \centering
  \begin{tikzpicture}[scale=1.0]
    % Círculo
    \draw[fill=verdecompacto!25,thick] (0,0) circle (0.8);
    \node at (0,-1.4) {\scriptsize Círculo};
    \node at (0,-1.8) {\scriptsize $G=0.886$};
    % Cuadrado
    \draw[fill=verdecompacto!40,thick] (2.6,-0.75) rectangle (4.1,0.75);
    \node at (3.35,-1.4) {\scriptsize Cuadrado};
    \node at (3.35,-1.8) {\scriptsize $G=1.000$};
    % Rectángulo moderado
    \draw[fill=naranjaextruido!25,thick] (5.4,-0.35) rectangle (8.0,0.35);
    \node at (6.7,-1.4) {\scriptsize Rect. moderado};
    \node at (6.7,-1.8) {\scriptsize $G=1.342$};
    % Rectángulo elongado
    \draw[fill=naranjaextruido!45,thick] (9.0,-0.18) rectangle (11.9,0.18);
    \node at (10.45,-1.4) {\scriptsize Rect. elongado};
    \node at (10.45,-1.8) {\scriptsize $G=1.775$};
  \end{tikzpicture}
  \vspace{0.8em}
  \begin{itemize}\footnotesize
    \item Valores apenas $>1.00$: desviaciones modestas de la forma eficiente.
    \item Valores $>1.342$: configuraciones claramente extruidas que restringen el desarrollo.
    \item En las regresiones se usa $\ln(\text{\Gindex})$ y se interpreta por \alert{cambios relativos / percentiles}, no por cambios unitarios.
  \end{itemize}
\end{frame}

%==============================================================================
\section{Datos}
%==============================================================================
\begin{frame}{Datos: el laboratorio luxemburgués}
  \begin{columns}[T]
    \begin{column}{0.58\textwidth}
      \begin{itemize}
        \item \textbf{4.297} transacciones de suelo urbano (arm's-length), \alert{2010--2014}.
        \item Cobertura nacional, mercado \alert{monocéntrico} (Luxemburgo-Ville concentra $>$50\% del empleo).
        \item $>$20 variables enriquecidas con SIG: físicas, locacionales, de vecindario y temporales.
        \item Polígonos catastrales (\textit{Plan Cadastral Numérisé}) $\rightarrow$ área $A_n$ y perímetro $P_n$.
        \item Variable dependiente: $\ln(\text{precio por área})$, en €/are (100 m\textsuperscript{2}).
      \end{itemize}
    \end{column}
    \begin{column}{0.42\textwidth}
      \centering
      \textbf{Tratamiento de outliers}
      \vspace{0.3em}
      \begin{table}\scriptsize
      \begin{tabular}{lrr}
        \toprule
        Año & Inicial & Limpio \\
        \midrule
        2010 & 949 & 820 \\
        2011 & 870 & 732 \\
        2012 & 859 & 769 \\
        2013 & 748 & 666 \\
        2014 & 871 & 759 \\
        \midrule
        \textbf{Total} & \textbf{4.297} & \textbf{3.746} \\
        \bottomrule
      \end{tabular}
      \end{table}
      \scriptsize 551 outliers removidos en 3 criterios (incl.\ outliers espaciales).
    \end{column}
  \end{columns}
\end{frame}

\begin{frame}{Estadísticas descriptivas (N = 3.746)}
  \centering
  \begin{table}\small
  \begin{tabular}{lrrrr}
    \toprule
    Variable & Media & D.E. & Mín & Máx \\
    \midrule
    \multicolumn{5}{l}{\itshape Dependiente}\\
    Precio por are (€) & 66.451 & 56.513 & 5.025 & 485.000 \\
    Precio por are (ln) & 10,835 & 0,745 & 8,52 & 13,09 \\
    \midrule
    \multicolumn{5}{l}{\itshape Características físicas}\\
    \Gindex (forma) & 1,30 & 0,32 & 1,08 & 9,63 \\
    Ln \Gindex & 0,25 & 0,16 & 0,07 & 2,27 \\
    Tamaño (are) & 6,48 & 9,78 & 0,01 & 269,23 \\
    \midrule
    \multicolumn{5}{l}{\itshape Selección de controles}\\
    Tiempo en auto al CBD (min) & 24,97 & 14,42 & 2,07 & 73,10 \\
    Zona altamente edificable & 0,89 & 0,31 & 0 & 1 \\
    \bottomrule
  \end{tabular}
  \end{table}
  \vspace{0.4em}
  \footnotesize El \Gindex medio es \alert{1,30} con amplio rango (1,08 a 9,63): variabilidad útil para identificar el efecto forma.
\end{frame}

%==============================================================================
\section{Metodología}
%==============================================================================
\begin{frame}{Estrategia empírica: modelo hedónico}
  Partimos del equilibrio de precios hedónicos de Rosen (1974):
  \[ p_{n,t}=f(C_{n,t};\,L_{n,t};\,V_{n,t};\,t) \]
  \footnotesize ($C$: físicas, $L$: locacionales, $V$: vecindario, $t$: tiempo)
  \vspace{0.5em}
  \normalsize
  \textbf{Modelo base log-lineal (OLS):}
  \[ \ln p_{n,t}=\alpha_0+\sum_{k=1}^{K}\beta_k^{\,t}\,x_{nk}^{t}+\delta_0^{\,t}+\varepsilon_n^{t} \]
  \begin{itemize}\small
    \item Forma entra como $\ln(\text{\Gindex})$ + términos no lineales e interacciones con el tamaño.
    \item Efectos fijos por \alert{municipio} (106 comunas) y por \alert{año}.
  \end{itemize}
\end{frame}

\begin{frame}{Dependencia espacial}
  \begin{itemize}
    \item Los precios del suelo son \alert{inherentemente espaciales}: ignorar la dependencia sesga e ineficiencia las estimaciones OLS.
    \item Diagnóstico: \textbf{Moran's I = 0,213} ($p=0.0000$) $\Rightarrow$ autocorrelación espacial positiva.
  \end{itemize}
  \vspace{0.3em}
  \textbf{Modelo de rezago espacial (SLM)}, estimado por GMM con corrección heterocedástica:
  \[ \ln p_{n,t}=\alpha_0+\lambda\!\!\sum_{j\neq n}\!w_{n,j}^{t}\ln p_{j,t}+\sum_{k}\beta_k^{t}x_{nk}^{t}+\delta_0^{t}+\varepsilon_n^{t} \]
  \begin{block}{Selección de modelo (estrategia GETS)}
    \footnotesize Desde el SARAR general $\rightarrow$ el coeficiente $\hat\rho$ resulta no significativo $\Rightarrow$ el modelo preferido se reduce al \alert{SLM}. Pesos espaciales: \textbf{$k$ vecinos más cercanos} ($k=10$).
  \end{block}
\end{frame}

\begin{frame}{Endogeneidad: ¿es la forma exógena?}
  Tres argumentos respaldan tratar la forma como \alert{predeterminada}:
  \vspace{0.4em}
  \begin{enumerate}
    \item \textbf{Predeterminación}: los límites catastrales son registros administrativos rígidos; reparcelar es costoso y raro ante una transacción puntual.
    \item \textbf{Ausencia de causalidad inversa}: la variación de forma proviene de subdivisiones históricas y restricciones topográficas, no de precios esperados.
    \item \textbf{Variables omitidas correlacionadas}: se controlan con (i) distancias continuas, (ii) características de vecindario, (iii) EF municipales y (iv) el rezago espacial del SLM.
  \end{enumerate}
  \vspace{0.4em}
  \footnotesize Un enfoque IV exigiría instrumentos difíciles de validar (fecha de formación catastral, régimen de planificación) $\rightarrow$ probablemente introduciría más sesgo del que corrige.
\end{frame}

%==============================================================================
\section{Resultados}
%==============================================================================
\begin{frame}{Resultado principal}
  \begin{block}{Hallazgo titular}
    \centering
    \large Una parcela con \Gindex $=1{,}342$ (moderadamente elongada) se transa con un \\[0.3em]
    \alert{\Huge $-11{,}3\%$} \\[0.3em]
    \large de descuento frente a la parcela de forma media (\Gindex $=1{,}3$), \\ manteniendo todo lo demás constante.
  \end{block}
  \vspace{0.4em}
  \begin{itemize}\small
    \item Las parcelas más extruidas se transan a precios significativamente menores que las compactas de igual tamaño y localización.
    \item El \Gindex se capitaliza \alert{sistemáticamente} en el valor del suelo.
    \item La forma es un atributo \alert{secundario pero económicamente relevante} (los derechos de desarrollo / zonificación siguen siendo de primer orden).
  \end{itemize}
\end{frame}

\begin{frame}{Resultados de estimación}
  \centering
  \begin{table}\scriptsize
  \begin{tabular}{lccc}
    \toprule
    Var.\ dep.: ln precio por are & (1) OLS & (2) OLS int. & (3) SLM-GMM \\
    \midrule
    ln \Gindex & $-0{,}260^{***}$ & $0{,}141$ & $0{,}068$ \\
    ln \Gindex$^2$ & $0{,}437^{***}$ & $-0{,}771^{***}$ & $-0{,}712^{***}$ \\
    ln tamaño (are) & $-0{,}163^{***}$ & $0{,}075^{*}$ & $0{,}074^{*}$ \\
    ln tamaño$^2$ & $-0{,}036^{***}$ & $-0{,}066^{***}$ & $-0{,}064^{***}$ \\
    ln \Gindex $\times$ ln tamaño & & $-0{,}431^{***}$ & $-0{,}391^{***}$ \\
    ln \Gindex$^2 \times$ ln tamaño$^2$ & & $0{,}041^{*}$ & $0{,}043^{***}$ \\
    \midrule
    Tiempo al CBD (min) & $-0{,}020^{***}$ & $-0{,}020^{***}$ & $-0{,}014^{***}$ \\
    Zona altamente edificable & $1{,}503^{***}$ & $0{,}893^{***}$ & $0{,}801^{***}$ \\
    Baulücke & $-0{,}056^{**}$ & $-0{,}071^{***}$ & $-0{,}073^{***}$ \\
    \midrule
    $\hat\lambda$ (rezago espacial) & & & $0{,}294^{***}$ \\
    EF municipio / año & Sí & Sí & Sí \\
    Pseudo-$R^2$ & 0,4149 & 0,4304 & 0,4210 \\
    \bottomrule
  \end{tabular}
  \end{table}
  \vspace{0.2em}
  \scriptsize N = 3.746. $^{*}p<0{,}1$; $^{**}p<0{,}05$; $^{***}p<0{,}01$. $\hat\lambda$ positivo y significativo: el entorno encarece la parcela.
\end{frame}

\begin{frame}{Lectura distribucional del efecto}
  \footnotesize El cambio unitario en el \Gindex es irrealista (de media a extremo). Comparamos cuatro parcelas de referencia (resto de variables en su media):
  \vspace{0.4em}
  \centering
  \begin{table}\small
  \begin{tabular}{lccr}
    \toprule
    Forma & \Gindex & Precio estimado & vs.\ media \\
    \midrule
    Cuadrado & 1,000 & € 319.104 & $+10{,}2\%$ \\
    \rowcolor{azulsecundario!15} Media muestral & 1,300 & € 289.617 & --- \\
    Moderadamente elongada & 1,342 & € 256.864 & \alert{$-11{,}3\%$} \\
    Altamente elongada (1:10) & 1,775 & € 186.323 & $-35{,}5\%$ \\
    \bottomrule
  \end{tabular}
  \end{table}
  \vspace{0.6em}
  \begin{itemize}\footnotesize
    \item El descuento del 11,3\% corresponde a un movimiento del \alert{rango intercuartílico} del \Gindex.
    \item El efecto es \alert{no lineal}: se acelera con mayor irregularidad.
    \item La interacción \Gindex$\times$tamaño es negativa: el castigo es \alert{mayor en parcelas grandes/valiosas}.
  \end{itemize}
\end{frame}

\begin{frame}{Heterogeneidad por tamaño de parcela}
  \begin{columns}[T]
    \begin{column}{0.5\textwidth}
      \centering
      \begin{tikzpicture}[scale=0.85]
        \draw[->] (0,0) -- (4,0) node[right]{\scriptsize ln \Gindex};
        \draw[->] (0,0) -- (0,3) node[above]{\scriptsize precio};
        % parcela pequeña: pendiente plana
        \draw[verdecompacto,very thick] (0.3,2.4) -- (3.6,2.0);
        \node[verdecompacto] at (2,2.6) {\scriptsize parcela pequeña};
        % parcela grande: pendiente pronunciada
        \draw[naranjaextruido,very thick] (0.3,2.4) -- (3.6,0.4);
        \node[naranjaextruido] at (2.4,0.9) {\scriptsize parcela grande};
      \end{tikzpicture}
    \end{column}
    \begin{column}{0.5\textwidth}
      \begin{itemize}\small
        \item \textbf{Parcelas pequeñas}: la irregularidad apenas reduce el valor (pendiente plana) — ya están geométricamente limitadas.
        \item \textbf{Parcelas grandes / valiosas}: la regularidad pesa mucho más (pendiente negativa pronunciada).
        \item Confirma el mecanismo: mayores costos / menor utilidad al desarrollar lotes irregulares cuando hay más valor en juego.
      \end{itemize}
    \end{column}
  \end{columns}
\end{frame}

%==============================================================================
\section{Robustez}
%==============================================================================
\begin{frame}{Verificaciones de robustez}
  El hallazgo del 11,3\% sobrevive a una batería de pruebas:
  \vspace{0.4em}
  \begin{columns}[T]
    \begin{column}{0.5\textwidth}
      \begin{itemize}\small
        \item \textbf{Medidas de forma alternativas}: Gravelius y Morton dan resultados cualitativamente idénticos.
        \item \textbf{Outliers}: el \Gindex sigue significativo con la muestra completa (cae el ajuste: pseudo-$R^2$ 0,42 $\to$ 0,37).
        \item \textbf{Controles socio-demográficos}: tenencia, \% extranjeros, ingresos — sin cambios materiales.
      \end{itemize}
    \end{column}
    \begin{column}{0.5\textwidth}
      \begin{itemize}\small
        \item \textbf{Submercados}: 5 regiones (cantones) — sin diferencias significativas.
        \item \textbf{Modelos espaciales}: OLS, SLM, SEM, SARAR — signos y significancia consistentes.
        \item \textbf{Estabilidad anual}: coeficientes consistentes año por año.
        \item \textbf{Diagnóstico visual}: el \Gindex no se concentra en zonas de bajo precio $\Rightarrow$ no es proxy de localización.
      \end{itemize}
    \end{column}
  \end{columns}
\end{frame}

\begin{frame}{Robustez: medidas de forma alternativas (SLM-GMM)}
  \centering
  \begin{table}\small
  \begin{tabular}{lccc}
    \toprule
    Var.\ dep.: ln precio por are & Gravelius & Morton & \Gindex \\
    \midrule
    ln forma & $0{,}068$ & $-0{,}029$ & $0{,}051$ \\
    ln forma$^2$ & $-0{,}712^{***}$ & $0{,}622^{***}$ & $-0{,}615^{**}$ \\
    ln tamaño (are) & $0{,}074^{*}$ & $0{,}074^{*}$ & $0{,}069^{*}$ \\
    ln forma $\times$ ln tamaño & $-0{,}391^{***}$ & $-0{,}276^{***}$ & $-0{,}381^{**}$ \\
    \midrule
    $\hat\lambda$ (rezago espacial) & $0{,}294^{***}$ & $0{,}291^{***}$ & $0{,}199^{***}$ \\
    N & 3.746 & 3.746 & 4.297 \\
    Pseudo-$R^2$ & 0,4210 & 0,4042 & 0,3716 \\
    \bottomrule
  \end{tabular}
  \end{table}
  \vspace{0.4em}
  \footnotesize Gravelius $\approx$ \Gindex (transformación lineal). Morton es empíricamente menos convincente. La última columna incluye outliers (N = 4.297) y muestra menor ajuste.
\end{frame}

%==============================================================================
\section{Conclusiones}
%==============================================================================
\begin{frame}{Contribuciones}
  \begin{block}{1. Contribución empírica (principal)}
    Primera cuantificación rigurosa del \alert{descuento por forma} en un mercado de suelo urbano. La parcela individual —unidad operativa del desarrollo— no había recibido atención empírica directa.
  \end{block}
  \begin{block}{2. Contribución metodológica (secundaria)}
    Introducción y validación del \alert{\Gindex}: normaliza la compacidad respecto a un \textbf{cuadrado} (forma eficiente), más interpretable para desarrolladores, tasadores y planificadores que la referencia circular.
  \end{block}
\end{frame}

\begin{frame}{Implicancias para la práctica}
  \begin{itemize}
    \item \textbf{Desarrolladores}: el descuento por forma es un factor de riesgo cuantificable en la adquisición, y una fuente de valor vía racionalización y ensamble de lotes.
    \item \textbf{Tasadores}: respaldo empírico para incluir un \alert{factor de ajuste por forma} en la valuación residual del suelo.
    \item \textbf{Planificadores y gobiernos locales}: las decisiones de subdivisión tienen consecuencias medibles sobre el mercado. Configuraciones compactas y ortogonales pueden generar mayor captura de valor.
  \end{itemize}
\end{frame}

\begin{frame}{Limitaciones y agenda futura}
  \begin{columns}[T]
    \begin{column}{0.5\textwidth}
      \textbf{Limitaciones}
      \begin{itemize}\small
        \item Luxemburgo: país pequeño, de altos ingresos y catastro rígido $\Rightarrow$ generalizar con cautela.
        \item Período 2010--2014: una fase específica del ciclo del mercado.
        \item Solo suelo urbano (se excluye agrícola/rural).
      \end{itemize}
    \end{column}
    \begin{column}{0.5\textwidth}
      \textbf{Investigación futura}
      \begin{itemize}\small
        \item Otros contextos urbanos (grilla americana vs.\ trazado europeo orgánico).
        \item Datos catastrales con variación temporal: ¿se intensifica el efecto con escasez de suelo?
        \item Combinar el \Gindex con datos 3D de forma edilicia: optimización conjunta parcela--edificio.
      \end{itemize}
    \end{column}
  \end{columns}
\end{frame}

\begin{frame}{Mensajes para llevar}
  \begin{enumerate}
    \item La \alert{forma de la parcela se capitaliza sistemáticamente} en el precio del suelo urbano.
    \item Pasar de una forma media a una moderadamente elongada implica un \alert{descuento de $\sim$11,3\%}.
    \item El efecto es \alert{no lineal} y \alert{más fuerte en parcelas grandes}.
    \item El \Gindex aporta una métrica \alert{validada e interpretable} para la economía urbana.
    \item Resultados \alert{robustos} a medidas alternativas, outliers, submercados y especificaciones espaciales.
  \end{enumerate}
\end{frame}

%------------------------------ Cierre ----------------------------------------
\begin{frame}[plain]
  \centering
  \vspace{1.5cm}
  {\Huge \textcolor{azulprincipal}{¡Gracias!}}\\[1cm]
  {\large \textbf{El efecto de la forma de la parcela sobre el precio del suelo urbano}}\\[0.5cm]
  {\normalsize Glumac \quad $\cdot$ \quad Herrera-Gómez \quad $\cdot$ \quad Licheron \quad $\cdot$ \quad Francke}\\[1cm]
  {\small Comentarios y preguntas bienvenidos}\\[0.3em]
  {\footnotesize \texttt{mherreragomez@conicet.gov.ar}}
\end{frame}

\end{document}
